\documentclass[12pt, twoside]{book}
\newcommand{\Przedmiot}{Informatyka śledcza}
%***************************************************************************************************
% Podstawowe ustawienia języka, według którego formatowany będzie dokument
\usepackage[polish]{babel}
\usepackage{diagbox}

% Pakiet babel dla polskiego języka powoduje konflikt z pakietem amssymb.
\let\lll\undefined

% W przypadku wielojęzykowości ustawia główny język dokumentu
\selectlanguage{polish}

% Kodowanie dokumentu
\usepackage[utf8]{inputenc} 
\usepackage[T1]{fontenc} 

% Dowolny rozmiar czcionek, kodowanie znaków
\usepackage{lmodern}

% Polskie wcięcia akapitów
\usepackage{indentfirst}

% Polskie formatowanie typograficzne
\frenchspacing

% Zapewnia liczne usprawnienia wyświetlania i organizacji matematycznych formuł. 
\usepackage{amsmath}
\usepackage{amssymb}
\usepackage{mathrsfs}
\usepackage{mathtools}

% ***************************************************************************************************
% Kolory  
% ***************************************************************************************************
% Umożliwia kolorowanie poszczególnych komórek tabeli i rozszerzoną kontrolę
\usepackage[table, svgnames]{xcolor}

% Umożliwia łatwą zmianę koloru linii w tabeli
\usepackage{tabu}

% Definicje kolorów
\definecolor{lgray}{HTML}{9F9F9F}
\definecolor{dgray}{HTML}{5F5F5F}

% ***************************************************************************************************
% Algorytmy i Listingi
% ***************************************************************************************************
\usepackage[ruled,vlined,linesnumbered,longend,algochapter]{algorithm2e}

% Zamiana nazwy środowiska z domyślnej "Algorithm X" na "Pseudokod X"
\newenvironment{algorithm-custom}[1][htb]{
	\renewcommand{\algorithmcfname}{Pseudokod}
	\begin{algorithm}[#1]%
	}{
\end{algorithm}
}

\newcommand\algcomment[1]{\footnotesize{#1}}
\SetCommentSty{algcomment}

% Wyśrodkowana tylda
\usepackage{textcomp}
\newcommand{\textapprox}{\raisebox{0.5ex}{\texttildelow}}

% Listowanie algorytmów
\renewcommand\algorithmcfname{Algorytm}
\renewcommand\listalgorithmcfname{Spis algorytmów}

% Pakiet listings i pełna polska konfiguracja (kodowanie UTF-8 oraz polskie nazwy spisów)
\usepackage{listings}
\renewcommand{\lstlistingname}{Listing}
\renewcommand{\lstlistlistingname}{Spis listingów}

\addto\captionspolish{%
	\renewcommand{\lstlistingname}{Listing}%
	\renewcommand{\lstlistlistingname}{Spis listingów}%
}

\lstset{
	inputencoding=utf8,
	extendedchars=true,
	literate=
		{ą}{{\k{a}}}1 {ć}{{\'c}}1 {ę}{{\k{e}}}1 {ł}{{\l{}}}1 {ń}{{\'n}}1
		{ó}{{\'o}}1 {ś}{{\'s}}1 {ź}{{\'z}}1 {ż}{{\.z}}1
		{Ą}{{\k{A}}}1 {Ć}{{\'C}}1 {Ę}{{\k{E}}}1 {Ł}{{\L{}}}1 {Ń}{{\'N}}1
		{Ó}{{\'O}}1 {Ś}{{\'S}}1 {Ź}{{\'Z}}1 {Ż}{{\.Z}}1
}

% ***************************************************************************************************
% Marginesy 
% ***************************************************************************************************
\usepackage[headheight=18pt, top=25mm, bottom=25mm, left=25mm, right=25mm]{geometry}

\makeatletter
\setlength\@fptop{0\p@}	
\makeatother

% ***************************************************************************************************
% Styl 
% ***************************************************************************************************
\usepackage{titling}
\usepackage[subfigure]{tocloft}	
\tocloftpagestyle{tableOfContentStyle}

\usepackage{fancyhdr}

% Domyślny czysty styl (używany w frontmatter - bez przedmiotu, bez numeru)
\fancypagestyle{plain}{
	\fancyhf{}
	\renewcommand{\headrulewidth}{0pt}
	\renewcommand{\footrulewidth}{0pt}
}

% Domyślny styl dla stron właściwych (środek rozdziału)
\fancypagestyle{custom}{
	\fancyhf{}
	\fancyhead[RO]{
		\hrulefill \hspace{16pt} \large Rozdział \thechapter
		\put(-471.5,5.5){\makebox(0,0)[l]{\small Politechnika Wrocławska}}
	}
	\fancyhead[LE]{
		\large Rozdział \thechapter \hspace{16pt} \hrulefill 
		\put(-188,5.5){\makebox(0,0)[l]{\small Wydział Informatyki i Telekomunikacji}}
	}
	\fancyfoot[LE,RO]{\thepage}
	\fancyfoot[LO,RE]{\small\Przedmiot}
	\renewcommand{\headrulewidth}{0pt}
	\renewcommand{\footrulewidth}{0.2pt}
}

% Styl dla List of Figures (BEZ PRZEDMIOTU)
\fancypagestyle{ListofFiguresStyle}{
	\fancyhf{}
	\fancyhead[RO]{
		\hrulefill \hspace{16pt} \large Spis rysunków
		\put(-471.5,5.5){\makebox(0,0)[l]{\small Politechnika Wrocławska}}
	}
	\fancyhead[LE]{
		\large Spis rysunków \hspace{16pt} \hrulefill 
		\put(-188,5.5){\makebox(0,0)[l]{\small Wydział Informatyki i Telekomunikacji}}
	}
	\fancyfoot[LE,RO]{\thepage}
	\renewcommand{\headrulewidth}{0pt}
	\renewcommand{\footrulewidth}{0.2pt}
}

% Styl dla List of Tables (BEZ PRZEDMIOTU)
\fancypagestyle{ListofTablesStyle}{
	\fancyhf{}
	\fancyhead[RO]{
		\hrulefill \hspace{16pt} \large Spis tabel
		\put(-471.5,5.5){\makebox(0,0)[l]{\small Politechnika Wrocławska}}
	}
	\fancyhead[LE]{
		\large Spis tabel \hspace{16pt} \hrulefill 
		\put(-188,5.5){\makebox(0,0)[l]{\small Wydział Informatyki i Telekomunikacji}}
	}
	\fancyfoot[LE,RO]{\thepage}
	\renewcommand{\headrulewidth}{0pt}
	\renewcommand{\footrulewidth}{0.2pt}
}

% Styl dla List of Listings
\fancypagestyle{ListofListingsStyle}{
	\fancyhf{}
	\fancyhead[RO]{
		\hrulefill \hspace{16pt} \large Spis listingów
		\put(-471.5,5.5){\makebox(0,0)[l]{\small Politechnika Wrocławska}}
	}
	\fancyhead[LE]{
		\large Spis listingów \hspace{16pt} \hrulefill 
		\put(-188,5.5){\makebox(0,0)[l]{\small Wydział Informatyki i Telekomunikacji}}
	}
	\fancyfoot[LE,RO]{\thepage}
	\renewcommand{\headrulewidth}{0pt}
	\renewcommand{\footrulewidth}{0.2pt}
}

% Styl dla List of Algorithms
\fancypagestyle{ListofAlgorithmsStyle}{
	\fancyhf{}
	\fancyhead[RO]{
		\hrulefill \hspace{16pt} \large Spis algorytmów
		\put(-471.5,5.5){\makebox(0,0)[l]{\small Politechnika Wrocławska}}
	}
	\fancyhead[LE]{
		\large Spis algorytmów \hspace{16pt} \hrulefill 
		\put(-188,5.5){\makebox(0,0)[l]{\small Wydział Informatyki i Telekomunikacji}}
	}
	\fancyfoot[LE,RO]{\thepage}
	\renewcommand{\headrulewidth}{0pt}
	\renewcommand{\footrulewidth}{0.2pt}
}

% Styl dla bibliografii (MA PRZEDMIOT)
\fancypagestyle{bibliographyStyle}{
	\fancyhf{}
	\fancyhead[RO]{
		\hrulefill \hspace{16pt} \large Bibliografia
		\put(-471.5,5.5){\makebox(0,0)[l]{\small Politechnika Wrocławska}}
	}
	\fancyhead[LE]{
		\large Bibliografia \hspace{16pt} \hrulefill 
		\put(-188,5.5){\makebox(0,0)[l]{\small Wydział Informatyki i Telekomunikacji}}
	}
	\fancyfoot[LE,RO]{\thepage}
	\fancyfoot[LO,RE]{\small\Przedmiot}
	\renewcommand{\headrulewidth}{0pt}
	\renewcommand{\footrulewidth}{0.2pt}
}

% Styl dla dodatków (MA PRZEDMIOT)
\fancypagestyle{appendixStyle}{
	\fancyhf{}
	\fancyhead[RO]{
		\hrulefill \hspace{16pt} \large Dodatki \thechapter
		\put(-472.1, 12.1){\makebox(0,0)[l]{\includegraphics[width=0.05\textwidth]{resources/pwr-logo}}}
		\put(-443,5.5){\makebox(0,0)[l]{\small Politechnika Wrocławska}}
	}
	\fancyhead[LE]{
		\large Dodatek \thechapter \hspace{16pt} \hrulefill 
		\put(-22, 12.1){\makebox(0,0)[l]{\includegraphics[width=0.05\textwidth]{wiz-logo}}}
		\put(-141,5.5){\makebox(0,0)[l]{\small Wydział Informatyki i Telekomunikacji}}
	}
	\fancyfoot[LE,RO]{\thepage}
	\fancyfoot[LO,RE]{\small\Przedmiot}
	\renewcommand{\headrulewidth}{0pt}
	\renewcommand{\footrulewidth}{0.2pt}
}

% Styl dla List of Abbreviation
\fancypagestyle{ListofAbbreviationStyle}{
	\fancyhf{}
	\fancyhead[RO]{
		\hrulefill \hspace{16pt} \large Spis skrótów
		\put(-471.5,5.5){\makebox(0,0)[l]{\small Politechnika Wrocławska}}
	}
	\fancyhead[LE]{
		\large Spis skrótów \hspace{16pt} \hrulefill 
		\put(-188,5.5){\makebox(0,0)[l]{\small Wydział Informatyki i Telekomunikacji}}
	}
	\fancyfoot[LE,RO]{\thepage}
	\renewcommand{\headrulewidth}{0pt}
	\renewcommand{\footrulewidth}{0.2pt}
}

% Styl dla pozostałych stron spisu treści (czysty)
\fancypagestyle{tableOfContentStyle}{
	\fancyhf{}
	\renewcommand{\headrulewidth}{0pt}
	\renewcommand{\footrulewidth}{0.2pt}
}

% Formatowanie tytułów rozdziałów i sekcji (wraz z subsubsection)
\usepackage{titlesec}
\setcounter{secnumdepth}{3}
\setcounter{tocdepth}{2}

\titleformat{\chapter}[hang]
{\vspace{-10ex}\normalfont\Huge\bfseries}{\thechapter.}{1ex}{} 

\titleformat{\section}[hang]
{\normalfont\Large\bfseries}{\thesection.}{1ex}{} 

\titleformat{\subsection}[hang]
{\normalfont\large\bfseries}{\thesubsection.}{1ex}{} 

\titleformat{\subsubsection}[hang]
{\normalfont\normalsize\bfseries}{\thesubsubsection.}{1ex}{}

% ***************************************************************************************************
% Linki
% ***************************************************************************************************
\usepackage{hyperref}
\hypersetup{
	colorlinks	=	true,
	linkcolor	=	blue,
	citecolor	=	blue,
	urlcolor	=	blue
}
\urlstyle{same}

% ***************************************************************************************************
% Listy
% ***************************************************************************************************
\usepackage{enumitem}
\newlist{myitemize}{itemize}{3}
\setlist[myitemize,1]{label=\textbullet, leftmargin=4mm}
\setlist[myitemize,2]{label=$\diamond$, leftmargin=8mm}
\setlist[myitemize,3]{label=$\diamond$, leftmargin=12mm}

% ***************************************************************************************************
% Inne pakiety
% ***************************************************************************************************
\usepackage{graphicx}
\usepackage{float}
\usepackage{caption}
\usepackage{footnote}
\usepackage{dirtree}
\usepackage{multirow}
\usepackage{calc}

% ***************************************************************************************************
% Matematyczne skróty
% ***************************************************************************************************
\newcommand{\RR}{\mathbb{R}}
\newcommand{\NN}{\mathbb{N}}
\newcommand{\QQ}{\mathbb{Q}}
\newcommand{\ZZ}{\mathbb{Z}}
\newcommand{\IMP}{\rightarrow}
\newcommand{\IFF}{\leftrightarrow}

% ***************************************************************************************************
% Środowiska
% ***************************************************************************************************
\newtheorem{theorem}{Twierdzenie}[chapter]
\newtheorem{lemma}{Lemat}[chapter]
\newtheorem{example}{Przykład}[chapter]
\newtheorem{corollary}{Wniosek}[chapter]
\newtheorem{definition}{Definicja}[chapter]

\newenvironment{proof}{\par\noindent \textbf{Dowód.}}{\begin{flushright}\vspace*{-6mm}\mbox{$\blacklozenge$}\end{flushright}}
\newenvironment{remark}{\bigskip \par\noindent \small \textbf{Uwaga.}}{\begin{small}\vspace*{4mm}\end{small}}

% Dodatkowe pakiety pomocnicze
\usepackage{psfrag}
\usepackage{amsfonts}
\usepackage{supertabular}
\usepackage{array}
\usepackage{tabularx}
\usepackage{hhline}
\usepackage{microtype}
\usepackage{booktabs}
\usepackage{makecell}
\usepackage{rotating}
\usepackage{multicol}
\usepackage{cuted}
\usepackage{colortbl}
\usepackage{adjustbox}
\usepackage{color,soul}
\usepackage{subfigure}
\usepackage{pdfpages}

% Wyłączenie generowania pustych stron (zastąpienie cleardoublepage zwykłym clearpage)
\renewcommand{\cleardoublepage}{\clearpage}

\usepackage{pifont}
\newcommand{\cmark}{\ding{51}}
\newcommand{\xmark}{\ding{55}}
\newcommand{\bftab}{\fontseries{b}\selectfont}

\newcolumntype{P}[1]{>{\raggedright\arraybackslash\noindent}p{#1}}
\newcolumntype{R}[1]{>{\raggedleft\arraybackslash}p{#1}}
\newcolumntype{L}[1]{>{\raggedright\arraybackslash}p{#1}}
\newcolumntype{C}[1]{>{\centering\arraybackslash}m{#1}}

\usepackage[nocompress]{cite}
\usepackage{svg}
\usepackage{tabto}
\usepackage{wrapfig}

\begin{document}
\frontmatter

% ==================== STRONA TYTUŁOWA ====================
\begin{titlepage}
\includepdf[pages=1]{images/TitlePage.pdf}
\end{titlepage}
\clearpage

% ==================== SPIS TREŚCI ====================
% W frontmatter styl 'plain' jest całkowicie czysty (brak numeru, brak przedmiotu)
\pagestyle{tableOfContentStyle}
\tableofcontents
\clearpage

% ==================== GŁÓWNA TREŚĆ (ROZDZIAŁY) ====================
\mainmatter
\pagestyle{custom}

% Dynamicznie definiujemy, że od tego momentu pierwsza strona KAŻDEGO rozdziału 
% ma zawierać ZARÓWNO numer strony, JAK I nazwę przedmiotu w stopce.
\fancypagestyle{plain}{
	\fancyhf{}
	\fancyfoot[LE,RO]{\thepage}
	\fancyfoot[LO,RE]{\small\Przedmiot}
	\renewcommand{\headrulewidth}{0pt}
	\renewcommand{\footrulewidth}{0.2pt}
}

\chapter{Wprowadzenie}\label{ch:wprowadzenie}

\section{Kontekst technologiczny}\label{sec:kontekst_technologiczny}

Dynamiczny rozwój technologii informatycznych oraz powszechna cyfryzacja procesów biznesowych i instytucjonalnych sprawiły, że cyberbezpieczeństwo stało się jednym z kluczowych filarów ciągłości działania współczesnych organizacji. Równolegle do postępu w obszarze mechanizmów defensywnych następuje ciągła ewolucja metod i wektorów ataków stosowanych przez ugrupowania cyberprzestępcze oraz zaawansowane grupy sponsorowane przez państwa (ang. \textit{Advanced Persistent Threat}, APT). Tradycyjny paradygmat ochrony stacji roboczych i serwerów, opierający się na klasycznych programach antywirusowych (ang. \textit{Endpoint Protection Platforms}, EPP) weryfikujących sygnatury plików statycznie zapisanych na wolumenie dyskowym, wykazuje obecnie daleko idącą nieskuteczność.

Współczesne rodziny złośliwego oprogramowania, w szczególności konie trojańskie zdalnego dostępu (ang. \textit{Remote Access Trojan}, RAT, np. Agent Tesla, AsyncRAT, Remcos) oraz wyspecjalizowane oprogramowanie kradnące tożsamość cyfrową i poświadczenia (ang. \textit{Infostealer}, np. RedLine, Vidar, LummaC2), wykorzystują złożone mechanizmy zaciemniania kodu (ang. \textit{obfuscation}) oraz unikania mechanizmów detekcji (ang. \textit{defense evasion}). Do fundamentalnych technik utrudniających statyczną analizę należą:
\begin{itemize}
    \item \textbf{Techniki bezplikowe oraz Living-off-the-Land (LotL / LOLBins):} wykorzystywanie wbudowanych narzędzi administracyjnych systemu operacyjnego (takich jak \code{PowerShell}, \code{WMI}, \code{certutil}, \code{mshta}) do pobierania, dekodowania i uruchamiania ładunków w sposób niewymagający zapisu binariów na nośniku pamięci masowej.
    \item \textbf{Wstrzykiwanie kodu do przestrzeni procesów:} zacieranie śladów poprzez alokację i egzekucję złośliwego kodu maszynowego w kontekście bezpieczeństwa zaufanych procesów systemowych (np. \file{explorer.exe}, \file{svchost.exe}) przy użyciu funkcji API jądra Windows (\code{VirtualAllocEx}, \code{WriteProcessMemory}, \code{CreateRemoteThread}).
    \item \textbf{Drążenie procesów (ang. \textit{Process Hollowing}):} technika polegająca na powołaniu legalnego procesu w stanie zawieszenia (\textit{suspended state}), odmapowaniu oryginalnej sekcji kodu (\code{NtUnmapViewOfSection}) i wstrzyknięciu w jej miejsce złośliwego pliku wykonywalnego przed wznowieniem wątku głównego.
    \item \textbf{Refleksywne ładowanie bibliotek (ang. \textit{Reflective DLL Injection}):} załadowanie biblioteki dynamicznej bezpośrednio z bufora pamięci RAM bez angażowania standardowego mechanizmu ładowania systemu (\code{ntdll!LdrLoadDll}) i bez rejestracji w oficjalnych strukturach modułów procesu (\code{PEB.Ldr}).
\end{itemize}

Powyższy stan rzeczy determinuje przesunięcie punktu ciężkości w operacjach bezpieczeństwa (ang. \textit{Security Operations Center}, SOC) oraz w informatyce śledczej i reagowaniu na incydenty (ang. \textit{Digital Forensics and Incident Response}, DFIR). Niezbędna staje się analiza ulotnych śladów binarnych w pamięci operacyjnej (RAM) oraz ciągła telemetria behawioralna jądra systemu operacyjnego. Ulotna pamięć RAM stanowi najbardziej wiarygodne źródło wiedzy o incydencie, gdyż przechowuje wstrzyknięte ładunki w postaci odszyfrowanej, struktury procesów, poświadczenia uwierzytelniające oraz aktywne deskryptory połączeń sieciowych.

\section{Problem badawczy i motywacja}\label{sec:problem_badawczy}

Współczesna inżynieria detekcji (ang. \textit{Detection Engineering}) zmaga się z dwoma krytycznymi wyzwaniami:
\begin{enumerate}
    \item \textbf{Zjawiskiem przeciążenia alertami (ang. \textit{alert fatigue}):} wynikającym z generowania nadmiernej liczby fałszywych alarmów (ang. \textit{False Positives}, FP) przez niedostatecznie zoptymalizowane lub zbyt ogólne reguły detekcyjne.
    \item \textbf{Niską odpornością wskaźników kompromitacji na modyfikacje:} zjawiskiem opisanym w koncepcji Piramidy Bólu Davida Bianco (ang. \textit{Pyramid of Pain}). Tradycyjne wskaźniki kompromitacji (ang. \textit{Indicators of Compromise}, IoC), takie jak skróty kryptograficzne plików (MD5, SHA256) czy adresy IP serwerów dowodzenia (ang. \textit{Command and Control}, C2), są trywialne do zneutralizowania przez adwersarza poprzez zwykłą rekompilację lub rotację infrastruktury.
\end{enumerate}

Aby skutecznie przeciwdziałać zaawansowanym zagrożeniom, konieczne jest tworzenie reguł detekcyjnych operujących na wyższych piętrach piramidy - wskaźnikach aktywności (ang. \textit{Indicators of Activity}, IoA) oraz technikach, taktykach i procedurach (ang. \textit{Tactics, Techniques, and Procedures}, TTP) skatalogowanych w matrycy MITRE ATT\&CK.

Wytworzenie precyzyjnych reguł wymaga jednak przeprowadzenia kontrolowanej infekcji w środowisku laboratoryjnym, pozyskania zrzutu pamięci operacyjnej oraz rejestracji szczegółowych dzienników zdarzeń (np. za pośrednictwem narzędzia Microsoft Sysmon). Następnie konieczne jest przeprowadzenie korelacji artefaktów z pamięci RAM (ujawnionych np. przez Volatility 3 lub MemProcFS) z telemetrią logów systemowych (parsowanych przy użyciu nowoczesnych narzędzi, takich jak Hayabusa czy Chainsaw).

Istotną luką w wielu opracowaniach laboratoryjnych jest brak spójnego, powtarzalnego podejścia, które w sposób zintegrowany łączy fazę inżynierii laboratoryjnej, poinfekcyjną analizę śledczą oraz bezpośrednią syntezę pozyskanych danych w postaci reguł YARA (dla pamięci RAM) oraz Sigma (dla dzienników zdarzeń). Niniejszy projekt stanowi odpowiedź na tę potrzebę.

\section{Cel i teza projektu}\label{sec:cele_projektu}

Głównym celem projektu jest \textbf{zaprojektowanie, wdrożenie i walidacja bezpiecznego, odizolowanego środowiska laboratoryjnego oraz opracowanie metodyki ekstrakcji i korelacji poinfekcyjnych artefaktów pamięci RAM i dzienników zdarzeń Windows w celu wytworzenia i weryfikacji wysokowiernych reguł detekcyjnych YARA i Sigma, a także skompilowania kompleksowej matrycy IoC/IoA.}

W ramach celu głównego wyznaczono następujące cele szczegółowe:
\begin{enumerate}
    \item \textbf{Opracowanie architektury i wdrożenie laboratorium badawczego:}
    Utworzenie odizolowanego segmentu wirtualnego, obejmującego stację ofiary (system Windows 10/11) ze specjalnie ukształtowanym profilem telemetrii Sysmon i dezaktywowanymi mechanizmami ochronnymi (Windows Defender, AMSI) oraz węzeł analityczny bazujący na systemie Linux (np. REMnux), wyposażony w usługi emulacji sieci (FakeNet-NG / INetSim).
    
    \item \textbf{Akwizycja i kontrolowana egzekucja materiału badawczego:}
    Bezpieczne pobranie z repozytoriów złośliwego oprogramowania (np. MalwareBazaar) próbek reprezentujących rodziny Infostealer oraz RAT, ich statyczna wstępna weryfikacja, detonacja w monitorowanym węźle oraz pozyskanie pełnego zrzutu pamięci operacyjnej z poziomu hiperwizora wraz z eksportem dzienników \file{.evtx}.
    
    \item \textbf{Informatyka śledcza pamięci ulotnej (Memory Forensics):}
    Identyfikacja ukrytych procesów, anomalii w strukturze drzewa procesów, zmapowanie nieprzypisanych pamięci VAD ze statusem \code{PAGE_EXECUTE_READWRITE} (detekcja iniekcji kodu) oraz ekstrakcja zdekodowanych ładunków binarnych bezpośrednio z RAM za pomocą narzędzi Volatility 3 oraz MemProcFS.
    
    \item \textbf{Analiza telemetrii systemowej i korelacja zdarzeń:}
    Zautomatyzowane przetworzenie logów systemowych za pomocą narzędzi Chainsaw oraz Hayabusa, rekonstrukcja łańcucha wykonania (drzewo procesów, parametry CLI, manipulacje rejestrem w celach persystencji, aktywność sieciowa) oraz powiązanie zachowań telemetrycznych z artefaktami pamięciowymi.
    
    \item \textbf{Inżynieria detekcji i weryfikacja empiryczna:}
    Zaprojektowanie autorskich reguł YARA ukierunkowanych na wykrywanie artefaktów pamięciowych w procesach oraz generycznych reguł Sigma wykrywających anomalie behawioralne w dziennikach systemowych, weryfikacja ich skuteczności oraz współczynnika fałszywych alarmów (FP) w czystym systemie referencyjnym, a także stworzenie matrycy wskaźników IoC oraz IoA.
\end{enumerate}

\section{Zakres projektu}\label{sec:zakres_projektu}

Zakres pracy obejmuje pięć głównych faz realizacyjnych, przedstawionych w tabeli~\ref{tab:fazy_projektu}.

\begin{table}[htbp]
\centering
\caption{Zestawienie faz realizacyjnych i komponentów projektu}
\label{tab:fazy_projektu}
\tablefont
\begin{tabularx}{\linewidth}{l L{4.2cm} Y}
\toprule
\hcell{Faza} & \hcell{Obszar prac} & \hcell{Główne narzędzia / technologie} \\
\midrule
I & Architektura i wdrożenie laboratorium & Hiperwizor (Type-2), izolowana sieć wewnętrzna (\textit{host-only}), FakeNet-NG, INetSim, PowerShell hardening/unhardening, Sysmon. \\
\addlinespace
II & Pozyskanie próbek i bezpieczna detonacja & Repozytorium MalwareBazaar, narzędzia statyczne (PEstudio, Floss), mechanizm snapshotów hiperwizora, akwizycja RAM i logów \file{.evtx}. \\
\addlinespace
III & Śledcza analiza pamięci RAM & Volatility 3, MemProcFS, wtyczki \code{malfind}, \code{pslist}, \code{netscan}, \code{vadinfo}, disasemblery (Ghidra, IDA). \\
\addlinespace
IV & Analiza i korelacja logów systemowych & Sysmon Modular, Chainsaw, Hayabusa, Windows Event Logs, parsery JSON/CSV, timeline analysis. \\
\addlinespace
V & Inżynieria detekcji i walidacja & Standard YARA (analiza pamięci procesu), standard Sigma (analityka logów), walidacja FP/TP, matryca IoC/IoA, MITRE ATT\&CK. \\
\bottomrule
\end{tabularx}
\end{table}

\subsection*{Ograniczenia i założenia projektowe}
W celu zachowania przejrzystości metodologicznej i powtarzalności eksperymentu przyjęto następujące założenia:
\begin{itemize}
    \item \textbf{Docelowy system operacyjny:} Eksperymenty przeprowadzane są w 64-bitowym środowisku Microsoft Windows (Windows 10/11), jako że stanowi ono dominujący cel w sektorze ataków na końcówki robocze.
    \item \textbf{Izolacja sieciowa:} Całkowity zakaz bezpośredniej komunikacji węzła infekcyjnego z publicznym Internetem; cała komunikacja C2 jest zamykana i emulowana lokalnie przez stację analityczną (FakeNet-NG / INetSim).
    \item \textbf{Faza ataku:} Badania koncentrują się na analizie poinfekcyjnej (\textit{post-execution / post-compromise}), kładąc główny nacisk na fazy wykonania (\textit{Execution}), utrzymania dostępu (\textit{Persistence}), unikania obrony (\textit{Defense Evasion}) oraz komunikacji (\textit{Command and Control}).
\end{itemize}

\section{Metodologia badawcza i warsztat technologiczny}\label{sec:metodologia}

Metodologia zastosowana w niniejszym projekcie opiera się na iteracyjnym cyklu inżynierii detekcji (ang. \textit{Detection Engineering Life Cycle}), który integruje analizę behawioralną, analizę statyczną artefaktów oraz formalne testowanie reguł detekcyjnych.

W toku realizacji projektu wykorzystywane są następujące grupy narzędzi:
\begin{itemize}
    \item \textbf{Środowisko wirtualizacji i symulacji sieci:} Wirtualne maszyny operujące w odciętej podsieci; węzeł analityczny przechwytuje zapytania DNS, protokoły HTTP/HTTPS, SMTP oraz surowe połączenia TCP, generując wiarygodne syntetyczne odpowiedzi serwera C2.
    \item \textbf{Telemetria systemowa (Sysmon):} Wykorzystanie rozbudowanego pliku konfiguracyjnego (m.in. na bazie standardu \textit{Sysmon-Modular}), pozwalającego na rejestrację m.in. powoływania wątków zdalnych (Event ID 8), uzyskiwania dostępu do procesów LSASS/innych (Event ID 10) oraz ładowania sterowników i bibliotek (Event ID 6 i 7).
    \item \textbf{Frameworki analizy pamięci (Volatility 3 / MemProcFS):} Zastosowanie technik eksploracji pamięci fizycznej i wirtualnej, omijających manipulacje adwersarza w tablicach jądra (\textit{DKOM - Direct Kernel Object Manipulation}) i pozwalających na ekstrakcję czystego kodu wykonywalnego z deskryptorów VAD.
    \item \textbf{Narzędzia szybkiego triage logów (Hayabusa / Chainsaw):} Przetwarzanie wielomegabajtowych plików \file{.evtx} z wykorzystaniem wbudowanych reguł Sigma w celu natychmiastowej rekonstrukcji osi czasu incydentu.
    \item \textbf{Formaty reguł detekcyjnych (YARA i Sigma):} Opracowanie reguł YARA w standardzie wersji 4.x ze szczególnym uwzględnieniem modułów \code{pe} i selekcji heksadecymalnej oraz reguł Sigma w formacie YAML zgodnych ze specyfikacją \textit{SigmaHQ}.
\end{itemize}


\clearpage

% ==================== BIBLIOGRAFIA ====================
\phantomsection
\pagestyle{bibliographyStyle}
\addcontentsline{toc}{chapter}{Bibliografia}

\nocite{*} 

\bibliographystyle{plabbrv-fixed}
\bibliography{bibliography}
\clearpage


% ==================== SPISY KOŃCOWE ====================
% Przywracamy tutaj styl 'plain' do wersji całkowicie czystej:
\fancypagestyle{plain}{
	\fancyhf{}
	\renewcommand{\headrulewidth}{0pt}
	\renewcommand{\footrulewidth}{0pt}
}

% Spis rysunków
\phantomsection
\pagestyle{ListofFiguresStyle}
\addcontentsline{toc}{chapter}{Spis rysunków}
\listoffigures
\clearpage

% Spis tabel
\phantomsection
\pagestyle{ListofTablesStyle}
\addcontentsline{toc}{chapter}{Spis tabel}
\listoftables
\clearpage

% Spis listingów
\phantomsection
\pagestyle{ListofListingsStyle}
\renewcommand{\lstlistlistingname}{Spis listingów}
\addcontentsline{toc}{chapter}{Spis listingów}
\lstlistoflistings
\clearpage

\end{document}